\section{功能设想}

CyberArm 定位为一台可编程的视觉交互桌面机械臂。系统包含四个关节运动轴和一路夹爪执行轴，可通过摄像头或姿态传感器获取输入，并在 ESP32-S3 上完成检测、状态管理和运动控制。除预置模式外，关节运动、目标检测和传感器读取等底层能力也通过串口开放给上位机，用户可以利用 Python SDK 重新组合动作与交互逻辑。

\subsection{四个核心模块}

系统按功能划分为以下四个模块：

\begin{enumerate}
  \item \textbf{运动执行}：五路舵机分别驱动底座、肩、肘、腕和夹爪。上位机既可以直接给定关节角，也可以给出末端目标坐标，由逆运动学模块计算可行的关节配置。
  \item \textbf{视觉感知}：摄像头采集桌面及机械臂前方区域的图像，根据当前模式执行人脸、指尖、颜色物体或手势检测。检测坐标经过指数移动平均（EMA）和死区处理后，再送入坐标映射与运动控制模块，以减少图像噪声引起的舵机抖动。
  \item \textbf{可编程接口}：自配置控制板引出 ESP32-S3 的原生 USB 接口，固件通过 USB 串口提供按行分隔的 JSON 指令。上位机可以发送运动、检测和模式切换请求；固件也可以在目标出现、丢失或动作完成时主动上报事件。
  \item \textbf{交互与反馈}：底座上的 OLED 用于显示当前模式、连接状态和故障信息。MPU6050 可作为独立姿态输入，在镜像模式下将测得的横滚和俯仰变化映射到机械臂的部分关节。
\end{enumerate}

\subsection{五种交互模式}

系统计划提供五种交互模式。其中前四种以摄像头为主要输入，第五种使用 MPU6050。模式切换由状态机统一管理，切换时会清空上一模式的滤波状态和未完成指令，避免旧数据继续驱动机械臂。

\textbf{人脸追随。}拟采用 esp32-camera 完成图像采集，并复用 ESP-WHO/ESP-DL 的轻量人脸检测组件。由于控制板和摄像头接口均自行配置，需要编写对应的板级引脚定义，并验证模型输入格式和帧缓冲设置。控制器根据人脸框中心与图像中心的偏差调整底座和俯仰关节，使机械臂朝向目标。连续若干帧未检测到人脸时，系统停止跟踪并以受限速度回到待机姿态。

\textbf{指尖追踪。}图像先转换到 YCbCr 色彩空间并进行肤色分割，再通过形态学处理、最大连通域和凸包特征估计指尖位置。首版只跟踪桌面平面内的位置变化，末端高度采用预设值。目标短时丢失时保持当前位置，超过设定时限后退出跟踪。

\textbf{物体跟踪。}用户在标定阶段选定目标颜色，系统在 HSV 空间内设定阈值，并以最大连通域的质心作为目标位置。该方案适合颜色单一、背景可控的桌面场景。对于短时遮挡，可在上一次位置附近设置有限搜索区域；若要完成抓取，还需预先给定抓取高度和夹爪动作序列。

\textbf{手势指令。}在完成肤色分割后，计算轮廓面积、凸包面积、凸包缺陷数和长宽比等几何特征，对少量预定义手势进行分类。识别结果需要连续多帧一致才会触发动作，以降低单帧误检造成的误动作。具体手势集合将在实测后确定，不把复杂手势识别作为首版目标。

\textbf{镜像复现。}MPU6050 固定在前臂或手腕处，通过有线 I²C 原型或外接无线节点向主控发送姿态数据。系统利用加速度计和陀螺仪融合估计横滚角与俯仰角，并按标定比例映射到指定关节。六轴 IMU 缺少绝对航向参考，偏航角只能短时积分并需要手动回中。单个 IMU 也无法恢复人体肩、肘各自的姿态，因此该模式实现的是有限自由度动作映射，而非严格的一比一复现。

\subsection{可编程是怎么实现的}

固件将移动、检测、传感器读取和动作播放等功能抽象为相对独立的操作，并通过 ESP32-S3 原生 USB 串口对外提供统一接口。PC 端 Python SDK 负责串口连接、消息编解码、超时处理和回调分发，使演示脚本可以在不修改固件的情况下组合这些能力。

这种划分主要解决三个实际问题：

\begin{itemize}
  \item 固件与演示逻辑可以分别开发，只要协议保持兼容，二者不必同步修改；
  \item 跟踪增益、动作顺序和模式组合等参数可以在上位机侧快速调整；
  \item 现场能够用简短脚本调用底层功能，直接说明接口是否完整、系统是否可复用。
\end{itemize}

SDK 的首版接口将控制在约 20 个方法以内，核心方法包括末端移动、关节控制、单次检测、模式切换、动作播放、状态查询和事件回调。

\subsection{系统的几种工作状态}

机械臂上电后进入待机状态，其他状态由串口命令或本地按键触发。状态切换时，控制任务会检查当前动作是否允许中断，并设置统一的速度、角度和超时限制。

\begin{itemize}
  \item \textbf{待机}：保持默认姿态，等待模式切换或运动指令；
  \item \textbf{视觉追踪}：运行人脸、指尖或颜色物体检测，并周期性更新目标位置；
  \item \textbf{手势响应}：识别预定义手势，确认后调用对应动作；
  \item \textbf{镜像模式}：读取 MPU6050 姿态并映射到指定关节；
  \item \textbf{脚本模式}：由上位机发送动作序列，固件负责限位检查、执行和状态回传。
\end{itemize}

\subsection{功能到实现的快速索引}

表~\ref{tab:function-index}列出了各项功能的主要输入、处理方法和输出。表中内容是当前方案，具体参数将在样机测试后调整。

\begin{table}[h]
\centering
\caption{功能与技术实现对照}
\label{tab:function-index}
\small
\begin{tabularx}{\linewidth}{l l Y Y}
\toprule
\textbf{功能} & \textbf{输入} & \textbf{算法/逻辑} & \textbf{输出}\\
\midrule
点到点运动 & 桌面坐标 $(x,y,z)$ & 逆运动学、限位检查、轨迹插值 & 四个关节运动轴协同动作\\
人脸追随 & OV2640 图像 & ESP-WHO 检测、EMA、死区 & 底座与俯仰关节调整\\
指尖追踪 & OV2640 图像 & YCbCr 分割、形态学处理、凸包 & 末端跟随桌面内指尖位置\\
物体跟踪 & OV2640 图像 & HSV 阈值、最大连通域质心 & 跟踪指定颜色物体\\
手势识别 & OV2640 图像 & 轮廓几何特征、多帧投票 & 手势触发预置动作\\
镜像复现 & MPU6050 姿态 & 姿态解算、角度标定与映射 & 有限自由度动作映射\\
动作库 & 预存关键帧序列 & 关节空间插值、限速 & 预置及自定义动作\\
PC 脚本 & 原生 USB 串口 JSON 消息 & Python SDK、事件回调 & 组合控制与状态记录\\
\bottomrule
\end{tabularx}
\end{table}
